\chapter{The Mean Behavior Poisson process}
\label{ch:mbpp}

\begin{motivation}
Here is the pivot of the book. We define the companion process whose deterministic
intensity is the expected Hawkes intensity, observe that it is therefore Poisson
(so its counts are independent and easy), and note that it shares parameters with
the Hawkes process we want. Everything afterwards is solving and exploiting one
equation.
\end{motivation}

\section{Definition and parameter correspondence}

\begin{definition}[MBPP]\label{def:ch4-mbpp}
Given a Hawkes process with exogenous $s$ and kernel $\phi$, the Mean Behavior
Poisson process is the inhomogeneous Poisson process whose intensity is the
expected Hawkes intensity,
\begin{equation}
\xi(t)=\E[\lambda(t)]=s(t)+\int_0^t\phi(t-u)\,\xi(u)\dd u .
\label{eq:ch4-mbpp}
\end{equation}
\end{definition}

The same two objects $(s,\phi)$ define both processes, so they are in
\emph{one-to-one parameter correspondence}: fit the MBPP, read off the Hawkes
parameters.

\begin{intuitionbox}
The MBPP is the Hawkes process with its randomness averaged out. Figure-2-style
experiments confirm it: the mean of thousands of simulated Hawkes intensities lies
right on top of the (deterministic) MBPP curve $\xi$.
\end{intuitionbox}

\section{Why it is Poisson --- and why that is the point}

$\xi$ is deterministic, hence so is the compensator $\Xi(t)=\int_0^t\xi$; by
Watanabe (Thm~\ref{thm:ch2-watanabe}) the MBPP is a genuine Poisson process. Its
bucket counts are therefore \emph{independent}
$\mathrm{Poisson}(\Xi(o_{i-1},o_i])$, and the joint law factorises:
\begin{equation}
\Prob\big(N(o_{i-1},o_i]=\sfont C_i,\ i=1,\dots,m\big)
=\prod_{i=1}^m \frac{\Xi(o_{i-1},o_i]^{\sfont C_i}}{\sfont C_i!}\,e^{-\Xi(o_{i-1},o_i]} .
\end{equation}

\begin{whybox}[: a deterministic mean buys a product likelihood]
The Hawkes process failed precisely because its compensator is random
(Ch.~\ref{ch:problem}). Replacing $\lambda$ by its mean $\xi$ makes the compensator
deterministic, which by Watanabe restores independent increments, which gives the
product above --- the interval-censored likelihood we will minimise
(Ch.~\ref{ch:fitting}). One substitution converts an intractable problem into a sum
of Poisson terms.
\end{whybox}

\section{In code}

The package's \code{MBPP} object holds a kernel and an exogenous function and
exposes the intensity and compensator (Ch.~\ref{ch:solving} explains how they are
computed):

\begin{lstlisting}[caption={Constructing and evaluating an MBPP.}]
from hawkes_calibration import MBPP, ExponentialKernel, Constant
import numpy as np

mbpp = MBPP(ExponentialKernel(kappa=0.6, theta=0.8), Constant(2.0, horizon=40))
t = np.linspace(0, 30, 600)
xi = mbpp.intensity(t)          # the deterministic mean intensity
Xi = mbpp.compensator(t)        # its cumulative integral (the Poisson means)
\end{lstlisting}

\begin{examplebox}[: the MBPP is the mean]
Experiment \code{exp5\_mbpp\_impulse.py} overlays $\xi$ on the Monte-Carlo mean of
thousands of Hawkes intensities for a single impulse and for rectangle, Dassios,
and sine exogenous functions; the closed-form $\xi$ matches the simulated mean to
Monte-Carlo error ($\sim$2--6\%). This is the empirical content of
Definition~\ref{def:ch4-mbpp}.
\end{examplebox}

\begin{pitfall}
The MBPP is \emph{not} the Hawkes process --- only its mean. It cannot reproduce the
clustering/over-dispersion of real Hawkes counts, which is why interval-censored
counts are over-dispersed relative to the MBPP's Poisson assumption
(Ch.~\ref{ch:identif}). The parameters transfer; the higher moments do not.
\end{pitfall}

\begin{recap}
The MBPP is the inhomogeneous Poisson process with intensity
$\xi=\E[\lambda]=s+\phi*\xi$. Being deterministic-mean it is Poisson (Watanabe), so
its counts are independent and its likelihood is a product; and it shares
parameters with the Hawkes process. Next: how to actually solve
\eqref{eq:ch4-mbpp}.
\end{recap}

\exercise{Argue that if $\phi=0$ (no excitation) the MBPP and the Hawkes process
coincide and both are the inhomogeneous Poisson process of rate $s$.}
\exercise{Why is the MBPP count distribution \emph{narrower} (less dispersed) than
the true Hawkes count distribution with the same parameters?}
